\documentclass[11pt,a4paper]{article}
\usepackage[utf8]{inputenc}
\usepackage[margin=1in]{geometry}
\usepackage{xcolor}
\usepackage{booktabs}
\usepackage{titlesec}
\usepackage{hyperref}

\definecolor{darkblue}{RGB}{0, 51, 102}
\hypersetup{colorlinks=true, linkcolor=darkblue, urlcolor=blue}

\titleformat{\section}{\normalfont\Large\bfseries\color{darkblue}}{\thesection.}{0.5em}{}
\titleformat{\subsection}{\normalfont\large\bfseries\color{darkblue}}{\thesubsection}{0.5em}{}

\title{\textbf{Hybrid Machine Learning Pipeline for Real-Time Malicious URL Classification}}
\author{}
\date{\today}

\begin{document}

\maketitle

\section{Goal}
The goal of this project is to build a machine learning pipeline that automatically classifies URLs as
 \textbf{benign, phishing, malware, or defacement}. The system works in two layers.
\begin{itemize}
    \item \textbf{Layer 1: Lexical analysis:} Inspecting the structure of the URL string itself (length, special characters, entropy, suspicious tokens).
    \item \textbf{Layer 2: Passive OSINT enrichment:} Querying external sources for domain age (WHOIS), hosting infrastructure (IPInfo), and DGA-like pattern detection.
\end{itemize}
The project was trained on a dataset of 600,000 labeled URLs covering all four classes (dataset found from kaggle).

\section{Pipeline}

\subsection{Lexical Features}
This feature was extracted directly from the URL string without external API calls. 
Features include URL length, domain length, number of dots/hyphens/digits, presence 
of IP address, @ symbol, double slash redirects, HTTPS token in domain, domain entropy,
and number of subdomains.

\subsection{Passive OSINT Features}
Enriched on a 2,000 URL sample due to API rate limits. Features include:
\begin{itemize}
    \item \texttt{domain\_age\_days}: Domain registration age from WHOIS.
    \item \texttt{asn\_org, country, region}: Hosting infrastructure metadata from IPInfo.
    \item \texttt{is\_dga\_like}: Flags domains with entropy $> 3.8$.
    \item \texttt{suspicious\_domain}: young domain and DGA-like entropy combined.
\end{itemize}

\newpage
\section{Results}

\begin{table}[h]
\centering
\begin{tabular}{l l r}
\toprule
\textbf{Model} & \textbf{Dataset} & \textbf{Macro F1} \\
\midrule
Random Forest & 600k lexical only & 0.9287 \\
XGBoost & 600k lexical only & 0.9251 \\
XGBoost (enriched) & 2k passive OSINT & 0.9140 \\
\bottomrule
\end{tabular}
\end{table}

The base Random Forest model trained on the full 600k dataset achieved 
the strongest performance with a macro F1 of 0.9287. On the other hand the passive OSINT model has the 
lowest F1 score.

\section{Limitations}
\begin{itemize}
    \item Passive OSINT enrichment was restricted to 2,000 samples due to API rate limits on free IPInfo and WHOIS tiers.
    \item WHOIS lookups fail on dead or taken-down domains, though a missing record is a useful signal in itself.
    \item VirusTotal enrichment was excluded due to the free tier limit of 4 requests per minute.
    \item Rare classes like malware and phishing have limited representation in the 2k enriched sample.
\end{itemize}

\section{Future Work}
\begin{itemize}
    \item Run passive OSINT enrichment on the full 600k dataset.
    \item Add VirusTotal reputation scores using a premium API key.
    \item Encode categorical OSINT features like country and ASN as model inputs.
    \item Deploy the saved model as a REST API for real-time classification.
\end{itemize}

\end{document}